% Part: history
% Chapter: biographies 
% Section: alan-turing 

\documentclass[../../../include/open-logic-section]{subfiles}

\begin{document}

\olfileid{his}{bio}{tur} 

\olsection{అలన్ ట్యూరింగ్}

అలన్ ట్యూరింగ్ 1912 జూన్ 23న లండన్‌లోని మైడా వేల్‌లో జన్మించారు. సిద్ధాంత కంప్యూటర్ విజ్ఞానశాస్త్ర పితామహుడిగా భావిస్తారు. భౌతిక విజ్ఞానశాస్త్రాలు, గణితంపై ట్యూరింగ్‌కు చిన్నతనంలోనే ఆసక్తి మొదలైంది. కానీ బాలుడిగా చదివిన పాఠశాలల్లో సాహిత్యం, ప్రాచీన గ్రంథాలకే ప్రాధాన్యం ఉండడంతో ఆయన ఆసక్తులకు తగిన చోటు దొరకలేదు. ఫలితంగా చదువులో ఆశించినంతగా రాణించలేదు; చాలామంది ఉపాధ్యాయుల మందలింపులు ఎదుర్కొన్నారు.

\olphoto{turing-alan}{అలన్ ట్యూరింగ్}

ట్యూరింగ్ కేంబ్రిడ్జ్‌లోని కింగ్స్ కళాశాలలో గణితం చదివారు. 1936లో గణించదగిన ప్రమేయం అనే భావనను కచ్చితంగా నిర్వచించడానికి, నిర్ణయ సమస్య అనిర్ణయనీయతను నిరూపించడానికి, ఇప్పుడు ట్యూరింగ్ యంత్రం అని పిలిచేదాన్ని రూపొందించారు. అయితే తన లాంబ్డా కలనశాస్త్రం ద్వారా ఆ ఫలితాన్ని ముందే నిరూపించిన అలోంజో చర్చ్ ఆయనకంటే ముందు వచ్చారు. అయినా చర్చ్ ఫలితాన్ని ప్రస్తావిస్తూ ట్యూరింగ్ వ్యాసం ప్రచురితమైంది. చర్చ్ ట్యూరింగ్‌ను ప్రిన్స్‌టన్‌కు ఆహ్వానించారు; అక్కడ ఆయన 1936--1938 మధ్య గడిపి, చర్చ్ పర్యవేక్షణలో డాక్టరేట్ పొందారు.

తర్కశాస్త్రంపై ఆసక్తి ఉన్నప్పటికీ భౌతిక విజ్ఞానశాస్త్రాలపై ట్యూరింగ్ తొలి ఆసక్తి కొనసాగింది. రెండవ ప్రపంచ యుద్ధంలో బ్లెచ్లీ పార్క్‌లోని బ్రిటిష్ సంకేత విశ్లేషణ విభాగంలో పనిచేసినప్పుడు ఆయన ఆచరణాత్మక నైపుణ్యాలు ఉపయోగపడ్డాయి. జర్మన్ నౌకాదళ సమాచారంలో వాడిన ఎనిగ్మా సంకేతాన్ని ఛేదించడంలో ట్యూరింగ్ కేంద్ర పాత్ర పోషించారు. గణాంకశాస్త్రం, గూఢలిపిశాస్త్రంలో ఆయన నైపుణ్యం, ఎలక్ట్రానిక్ యంత్రాల ప్రవేశంతో కలిసి, ``బాంబ్'' అనే సంకేత విప్పే యంత్రాన్ని నిర్మించి ఆ సంకేతాన్ని ఛేదించే సామర్థ్యాన్ని బృందానికి ఇచ్చాయి. బ్లెచ్లీ పార్క్‌లో జర్మన్ లోరెన్జ్ గూఢలిపిని ఛేదించడానికి వాడిన, ప్రపంచంలోని తొలి ప్రోగ్రామ్ చేయదగిన ఎలక్ట్రానిక్ కంప్యూటర్ కొలోసస్ రూపకల్పనకు కూడా ఆయన ఆలోచనలు తోడ్పడ్డాయి.

ట్యూరింగ్ స్వలింగ ఆకర్షణ కలిగినవారు. అయినప్పటికీ 1942లో బ్లెచ్లీ పార్క్‌లో తన సహచరురాలైన జోన్ క్లార్క్‌కు వివాహ ప్రతిపాదన చేశారు; తరువాత నిశ్చితార్థాన్ని విరమించి, తాను స్వలింగ ఆకర్షణ కలిగినవాడినని ఆమెకు చెప్పారు. జీవితంలో ఆయనకు పలువురు ప్రేమికులు ఉన్నారు; అయితే అప్పట్లో యునైటెడ్ కింగ్‌డమ్‌లో స్వలింగ సంబంధ చర్యలు నేరాలు. 1952లో ఆయన అప్పటి ప్రేమికుడి స్నేహితుడు ట్యూరింగ్ ఇంట్లో చోరీ చేశాడు. పోలీసులకు ఫిర్యాదు చేసినప్పుడు, స్వలింగ చర్యలను ప్రభుత్వం త్వరలో చట్టబద్ధం చేస్తుందని భావించి, తనకు స్వలింగ సంబంధం ఉందని ట్యూరింగ్ చెప్పారు. అది నిజం కాదు; ఆయనపై తీవ్రమైన అసభ్య ప్రవర్తన ఆరోపణ మోపారు. జైలుకు వెళ్లడానికి బదులు లైంగిక వాంఛను తగ్గించే హార్మోన్ చికిత్సను ఎంచుకున్నారు. 1954 జూన్ 8న సయనైడ్ మోతాదు అధికం కావడంతో ఆయన మృతదేహం కనిపించింది---అది ఆత్మహత్య అయ్యే అవకాశమే ఎక్కువ. 2013లో రెండవ ఎలిజబెత్ రాణి ఆయనకు రాజమన్నింపు ఇచ్చారు.

\begin{reading}
అలన్ ట్యూరింగ్ సమగ్ర జీవిత చరిత్రకు \citet{Hodges2014} చూడండి. ఆయన జీవితం, కృషి స్ఫూర్తితో వచ్చిన \emph{Breaking the Code} నాటకం 1996లో డెరెక్ జాకోబీ ట్యూరింగ్ పాత్రలో నటించిన టీవీ రూపంలో నిర్మితమైంది. అకాడమీ అవార్డుకు నామినేట్ అయిన, బెనెడిక్ట్ కంబర్‌బాచ్, కీరా నైట్లీ నటించిన \emph{The Imitation Game} చిత్రం కూడా ట్యూరింగ్ జీవితం, బ్లెచ్లీ పార్క్ కాలం ఆధారంగా కొంత స్వేచ్ఛగా రూపొందింది \citep{Imitation2014}.

\citet{Radiolab2012}లో ట్యూరింగ్ జీవితం, కృషిపై అనేక పాడ్‌కాస్ట్‌లు ఉన్నాయి. బీబీసీ హొరైజన్ డాక్యుమెంటరీ \emph{The Strange Life and Death of
  Dr.~Turing} ఆన్‌లైన్‌లో చూడవచ్చు \citep{Sykes1992}. ట్యూరింగ్ శతజయంతి గౌరవార్థం 2012లో తయారైన, పనిచేసే లెగో ట్యూరింగ్ యంత్రపు చిన్న వీడియో \citep{Theelen2012}.

ట్యూరింగ్ యంత్రాలు, నిర్ణయ సమస్యపై ట్యూరింగ్ మూల వ్యాసం \citet{Turing1937}.
\end{reading}

\end{document}
